\documentclass[11pt,a4paper]{article}
\usepackage{curso}
\begin{document}
\portada{Ejercicios --- Lección}{4.3}{MCMC: Metropolis--Hastings y Gibbs}
\section{Ejercicios}
\begin{enumerate}[label=\textbf{E\arabic*.}]
\item (B) Demuestre el balance detallado de Metropolis con propuesta simétrica.
\item (B) Derive $P(s_i=+1\mid s_{-i})=\sigma(2\beta u_i)$ desde la energía de Ising.
\item (B) ¿Cuántos colores necesita una red triangular? ¿Un grafo de grado 16?
\item (B) ¿Por qué falla la actualización totalmente paralela aunque cada condicional sea correcta?
\item (I) Muestre que Gibbs es Metropolis--Hastings con propuesta $q=p(x_i\mid x_{-i})$ y aceptación 1.
\item (I) Para dos espines acoplados ($J=1$, $\beta=1$), escriba la matriz de transición de la actualización síncrona y compruebe que su
estacionaria no es Boltzmann.
\item (A) Implemente los tres muestreadores de la lección y reproduzca la tabla de energías a $T=3$ y $T=1{,}8$.
\end{enumerate}
\section{Soluciones}
\begin{solucion}
\textbf{E1.} Si $E(s')>E(s)$: $\pi(s)\,q\,e^{-\beta\Delta E}=\pi(s')\,q\cdot1$ porque $\pi(s')/\pi(s)=e^{-\beta\Delta E}$; simétrico en el otro caso.

\textbf{E2.} Los términos con $s_i$ son $-s_iu_i$; $P(+1)/P(-1)=e^{2\beta u_i}$.

\textbf{E3.} Triangular: 3 colores. Grado 16: a lo sumo 17 (Brooks: 16 si no es completo ni ciclo impar).

\textbf{E4.} La actualización conjunta no es una composición de pasos de Gibbs individuales: usa estados viejos de variables que cambian a la vez.

\textbf{E5.} Con $q(s\to s')=p(s_i'\mid s_{-i})$, el cociente de MH vale exactamente 1.

\textbf{E6.} Estados $(++),(+-),(-+),(--)$; cada espín se actualiza con $\sigma(\pm2)$ según el valor \emph{viejo} del otro. La matriz
resultante es doblemente estocástica y su estacionaria es uniforme, $P=\tfrac14$ para cada estado, mientras que Boltzmann da
$P(++)=P(--)=0{,}440$ y $P(+-)=P(-+)=0{,}060$.

\textbf{E7.} $T=3$: $-0{,}82$, $-0{,}79$, $+0{,}02$. A $T=1{,}8$ desde el estado ordenado las tres coinciden ($\approx-1{,}85$): el fallo síncrono
aparece al desordenarse.
\end{solucion}
\end{document}
